\documentclass[11pt,a4paper]{article}
\usepackage{../common/td_style}

\begin{document}

\makeuniversityheader{Travaux Dirigés : Série 00}{Introduction aux Biostatistiques \& Rappels Fondamentaux}{\textcolor{BioNavy}{\textbf{SÉRIE D'EXERCICES}}}

\begin{exercicebox}{Exercice 1 : Typologie des Variables en Biochimie Clinique}
Dans le cadre d'un bilan métabolique réalisé sur une cohorte de 60 patients suivis pour stéatose hépatique, l'équipe hospitalière relève les paramètres suivants :
\begin{enumerate}
  \item Le sexe du patient (Masculin / Féminin).
  \item La concentration sérique en triglycérides (en g/L).
  \item Le stade de fibrose hépatique évalué par biopsie (F0 : absence, F1 : minime, F2 : modérée, F3 : sévère, F4 : cirrhose).
  \item Le nombre d'épisodes d'ictère enregistrés au cours des 5 dernières années.
  \item L'absorbance spectrophotométrique à $540\,\text{nm}$ lors du dosage colorimétrique de la bilirubine.
  \item La présence ou absence de la mutation polymorphique du gène \textit{PNPLA3} (Sauvage / Hétérozygote / Homozygote muté).
\end{enumerate}

\textbf{Questions :}
\begin{enumerate}[label=\textbf{\alph*.}]
  \item Pour chacune de ces 6 variables, précisez sa nature exacte : qualitative (nominale ou ordinale) ou quantitative (discrète ou continue).
  \item Précisez pour chaque variable l'indicateur de tendance centrale (Moyenne, Médiane, Mode) et de dispersion (Écart-type, Écart interquartile) le plus pertinent à calculer.
\end{enumerate}
\end{exercicebox}

\begin{exercicebox}{Exercice 2 : Variabilité Expérimentale \& Mesures de Dispersion (Albumine Sérique)}
Le dosage de la concentration sérique en albumine (en g/L) a été réalisé en sextuplicat ($N = 6$ réplicats biologiques indépendants) chez des rats témoins sains :
\begin{center}
  \textbf{Valeurs observées :} \quad $38.2, \quad 41.5, \quad 39.8, \quad 42.0, \quad 37.5, \quad 41.0 \quad (\text{en g/L})$
\end{center}

\textbf{Questions :}
\begin{enumerate}[label=\textbf{\alph*.}]
  \item Calculez la moyenne empirique $\bar{x}$, la variance corrigée $s^2$ et l'écart-type $s$ de cet échantillon.
  \item Déterminez le coefficient de variation ($CV = s / \bar{x} \times 100$). Ce dosage présente-t-il une bonne répétabilité analytique (critère usuel : $CV < 5\%$) ?
  \item Expliquez la différence fondamentale entre la variance brute ($N$ au dénominateur) et la variance d'échantillon sans biais ($N-1$ au dénominateur).
\end{enumerate}
\end{exercicebox}

\begin{exercicebox}{Exercice 3 : SD versus SEM \& Intervalle de Confiance à 95\%}
Une équipe étudie l'activité spécifique de la catalase érythrocytaire (en U/mg de protéine) chez des sujets exposés à un polluant environnemental.
\begin{itemize}
  \item Sur un groupe de $N = 25$ patients, la moyenne observée est $\bar{x} = 142.0\,\text{U/mg}$ avec un écart-type $s = 20.0\,\text{U/mg}$.
\end{itemize}

\textbf{Questions :}
\begin{enumerate}[label=\textbf{\alph*.}]
  \item Calculez l'erreur standard de la moyenne ($SEM = s / \sqrt{N}$).
  \item Quelle est la différence biologique majeure entre rapporter un résultat sous la forme $\bar{x} \pm SD$ versus $\bar{x} \pm SEM$ ? Lequel reflète la variabilité biologique des patients ?
  \item En supposant la distribution normale, calculez l'Intervalle de Confiance à 95\% (IC 95\%) de la moyenne vraie de la population ($\mu$) sachant que $t_{0.05, \, 24} = 2.064$.
  \item Si l'équipe avait mesuré $N = 100$ patients avec la même dispersion ($s=20.0$), quelle aurait été la nouvelle valeur du $SEM$ et comment évoluerait la largeur de l'IC 95\% ?
\end{enumerate}
\end{exercicebox}

\begin{exercicebox}{Exercice 4 : Logique du Test d'Hypothèse \& Démystification de la $p$-value}
Un extrait aqueux de plante médicinale est testé pour son potentiel hypoglycémiant chez le rat rendu diabétique.
\begin{itemize}
  \item La glycémie moyenne du groupe non traité (Témoin, $N_1 = 10$) est de $3.20\,\text{g/L}$ ($s_1 = 0.40$).
  \item Le groupe recevant l'extrait ($N_2 = 10$) présente une glycémie moyenne de $2.65\,\text{g/L}$ ($s_2 = 0.35$).
  \item Le test $t$ de Student bilatéral fournit une $p$-value de $p = 0.0042$.
\end{itemize}

\textbf{Questions :}
\begin{enumerate}[label=\textbf{\alph*.}]
  \item Formulez clairement l'hypothèse nulle ($H_0$) et l'hypothèse alternative ($H_1$).
  \item Que signifie précisément l'affirmation « $p = 0.0042$ » ? Donnez sa définition probabiliste rigoureuse.
  \item Est-il exact de dire que « l'extrait a 99.58\% de chances d'être efficace » ? Expliquez pourquoi cette interprétation est erronée.
  \item Au seuil de risque $\alpha = 0.05$, quelle conclusion biologique et statistique doit-on formuler ?
\end{enumerate}
\end{exercicebox}

\begin{exercicebox}{Exercice 5 : Contrôle de Normalité \& Détection d'Outliers}
Lors d'une cinétique d'hémolyse induite par un détergent, les pourcentages d'hémolyse mesurés sur 8 tubes préparés identiquement sont :
\begin{center}
  $22.1, \quad 23.4, \quad 21.8, \quad 22.9, \quad 48.7, \quad 24.1, \quad 23.0, \quad 22.5 \quad (\%)$
\end{center}

\textbf{Questions :}
\begin{enumerate}[label=\textbf{\alph*.}]
  \item Déterminez la médiane, le premier quartile ($Q_1$), le troisième quartile ($Q_3$) et l'écart interquartile ($IQR = Q_3 - Q_1$).
  \item Appliquez la règle de Tukey pour identifier les valeurs aberrantes (outliers) : $[Q_1 - 1.5 \cdot IQR, \, Q_3 + 1.5 \cdot IQR]$. La valeur $48.7\%$ est-elle un outlier ?
  \item Quelle attitude le biochimiste doit-il adopter face à cette valeur aberrante avant de décider de la conserver ou de l'exclure ?
  \item Si l'échantillon ne suit pas une loi normale, quel test d'hypothèse non-paramétrique doit remplacer le test $t$ de Student pour comparer deux groupes ?
\end{enumerate}
\end{exercicebox}

\begin{exercicebox}{Exercice 6 : Théorème de Bayes \& Performances Diagnostiques d'un Test Sérologique (ELISA)}
Dans le cadre d'un programme de dépistage sérologique d'une hépatite virale chronique au sein d'une population générale asymptomatique :
\begin{itemize}
  \item La prévalence de l'infection dans cette population est estimée à $P(M) = 1\%$ ($0.01$).
  \item La sensibilité clinique du test immuno-enzymatique ELISA est de $Se = P(+ \mid M) = 98\%$ ($0.98$).
  \item La spécificité clinique du test ELISA est de $Sp = P(- \mid \bar{M}) = 95\%$ ($0.95$).
\end{itemize}

\textbf{Questions :}
\begin{enumerate}[label=\textbf{\alph*.}]
  \item Donnez la probabilité qu'un sujet sain soit testé faussement positif : $P(+ \mid \bar{M})$.
  \item Calculez la probabilité totale qu'un individu prélevé au hasard dans la population obtienne un résultat positif au test ELISA : $P(+)$ (formule des probabilités totales).
  \item En appliquant le Théorème de Bayes, calculez la \textbf{Valeur Prédictive Positive ($VPP = P(M \mid +)$)}, c'est-à-dire la probabilité qu'un individu testé positif soit réellement porteur du virus.
  \item Commentez la valeur obtenue. Pourquoi observe-t-on un tel « paradoxe des faux positifs » ($> 80\%$ de faux positifs parmi les tests positifs) malgré une sensibilité et une spécificité très élevées ?
  \item Calculez la \textbf{Valeur Prédictive Négative ($VPN = P(\bar{M} \mid -)$)}. Quelle est son utilité pratique pour exclure la maladie ?
  \item Si ce même test est désormais utilisé en milieu hospitalier spécialisé chez des patients présentant des symptômes hépatiques avérés (prévalence $P(M) = 30\%$), quelle devient la nouvelle valeur de la $VPP$ ? Que concluez-vous sur la dépendance de la $VPP$ vis-à-vis de la prévalence ?
\end{enumerate}
\end{exercicebox}

\begin{exercicebox}{Exercice 7 : Lois de Probabilités en Microbiologie \& Génétique (Poisson \& Binomiale)}
\textbf{Partie 1 : Contrôle de Stérilité d'un Bioréacteur (Loi de Poisson)} \\
Lors du suivi microbiologique d'un fermenteur industriel, on étale $100\,\mu\text{L}$ de prélèvement sur boîte de gélose. Le nombre de colonies bactériennes contaminantes (UFC) par boîte suit une loi de Poisson de paramètre $\lambda = 3.0$ colonies/boîte.
\begin{enumerate}[label=\textbf{\alph*.}]
  \item Rappelez les conditions biologiques justifiant la modélisation du comptage de colonies par une loi de Poisson $\mathcal{P}(\lambda)$.
  \item Calculez la probabilité que la boîte soit parfaitement stérile (aucune colonie observée, $X = 0$).
  \item Calculez la probabilité d'observer exactement 2 colonies ($X = 2$).
  \item Quelle est la probabilité d'observer une contamination notable, c'est-à-dire au moins 4 colonies ($X \ge 4$) ?
\end{enumerate}

\vspace{2mm}
\textbf{Partie 2 : Sélection de Mutants Résistants aux Antibiotiques (Loi Binomiale)} \\
Une culture bactérienne est soumise à un rayonnement UV pour induire des mutations ciblées. La probabilité qu'un clone bactérien acquière une résistance à la kanamycine est $p = 0.05$. On sélectionne au hasard $n = 10$ clones indépendants. Soit $Y$ la variable aléatoire désignant le nombre de clones résistants parmi les 10.
\begin{enumerate}[label=\textbf{\alph*.}]
  \setcounter{enumi}{4}
  \item Justifiez pourquoi $Y$ suit une loi binomiale $\mathcal{B}(n, p)$ et précisez son espérance mathématique $\mathbb{E}(Y)$ et son écart-type $\sigma_Y$.
  \item Calculez la probabilité d'obtenir exactement un clone résistant ($Y = 1$).
  \item Calculez la probabilité d'obtenir au moins un clone résistant ($Y \ge 1$). Combien de clones faudrait-il cribler pour être sûr à 95\% d'isoler au moins un mutant résistant ?
\end{enumerate}
\end{exercicebox}

\begin{exercicebox}{Exercice 8 : Loi Normale, $Z$-Score \& Intervalles de Référence Cliniques}
Dans une population de référence d'adultes sains, la glycémie à jeun suit une distribution normale de paramètres $\mu = 0.95\,\text{g/L}$ et $\sigma = 0.10\,\text{g/L}$ : $X \sim \mathcal{N}(0.95, \, 0.10^2)$.
\begin{enumerate}[label=\textbf{\alph*.}]
  \item Calculez l'intervalle de référence physiologique contenant 95\% des individus sains selon la règle canonique $[\mu \pm 1.96\sigma]$.
  \item Un patient se présente au laboratoire avec une glycémie mesurée à $1.22\,\text{g/L}$. Calculez son $Z$-score (valeur centrée-réduite $Z = \frac{X - \mu}{\sigma}$).
  \item Sachant que pour la loi normale standard, $\Phi(2.70) = P(Z \le 2.70) = 0.9965$, calculez le pourcentage d'individus sains ayant une glycémie supérieure ou égale à celle de ce patient. Cette valeur plaide-t-elle pour un état pathologique ?
  \item En toxicologie hépatique, l'activité sérique de l'alanine aminotransférase (ALAT) présente une distribution fortement asymétrique étirée vers les valeurs élevées (présence de patients cytolytiques). Pourquoi l'intervalle $\mu \pm 2\sigma$ est-il rigoureusement inapproprié dans ce cas ? Quelle transformation mathématique applique-t-on en routine pour normaliser ces données avant analyse ?
\end{enumerate}
\end{exercicebox}

\begin{exercicebox}{Exercice 9 : Puissance Statistique ($1-\beta$), Risque de 2\textsuperscript{de} Espèce ($\beta$) \& Éthique 3R}
Une équipe de biochimie pharmacologique étudie un peptide candidat pour réduire la taille d'une plaque d'athérome sur modèle murin.
\begin{enumerate}[label=\textbf{\alph*.}]
  \item Définissez rigoureusement le risque de première espèce $\alpha$ et le risque de seconde espèce $\beta$. Lequel correspond à un « faux positif » et lequel à un « faux négatif » ?
  \item Qu'appelle-t-on la \textbf{puissance statistique ($1 - \beta$)} d'une expérience ? Quel est le seuil standard de puissance exigé par les revues internationales et les agences du médicament (ex: FDA, EMA) ?
  \item Par contrainte budgétaire, le chercheur décide d'utiliser seulement $N = 3$ souris par lot. Un calcul de puissance rétrospectif révèle une puissance de seulement $25\%$ ($1 - \beta = 0.25$).
  \begin{itemize}
    \item Quelle est la probabilité que cette étude passe à côté d'un effet thérapeutique pourtant bien réel ?
    \item En quoi cette démarche viole-t-elle gravement le principe éthique des \textbf{3R (Remplacer, Réduire, Raffiner)} ? Pourquoi « sacrifier trop peu d'animaux » est-il tout aussi contraire à l'éthique que « sacrifier trop d'animaux » ?
  \end{itemize}
  \item Citez trois stratégies méthodologiques concrètes permettant d'accroître la puissance statistique d'un protocole sans nécessairement augmenter le nombre total d'animaux.
\end{enumerate}
\end{exercicebox}

\end{document}
